\documentclass[11pt,a4paper]{article}

% ---------- Langue / encodage ----------
\usepackage[T1]{fontenc}
\usepackage[utf8]{inputenc}
\usepackage[french]{babel}
\usepackage{lmodern}
\usepackage{microtype}

\usepackage{tabularx}
\usepackage{array}

% ---------- Mise en page ----------
\usepackage[a4paper,margin=2.2cm]{geometry}

% ---------- Maths / symboles ----------
\usepackage{amsmath,amssymb}
\usepackage{pifont} % \ding{51} \ding{55}
\newcommand{\cmark}{\ding{51}}
\newcommand{\xmark}{\ding{55}}

% ---------- Couleurs ----------
\usepackage{xcolor}
\definecolor{sectioncolor}{HTML}{1A5276}    % bleu foncé — sections
\definecolor{subsectioncolor}{HTML}{2E86C1} % bleu moyen — subsections
\definecolor{questioncolor}{HTML}{6C3483}   % violet — questions
\definecolor{codebg}{HTML}{F7F7F7}          % gris très clair — fond code
\definecolor{codeframe}{HTML}{CCCCCC}       % gris — cadre code
\definecolor{keyword}{HTML}{0077AA}         % bleu — mots-clés
\definecolor{string}{HTML}{2E7D32}          % vert — chaînes
\definecolor{comment}{HTML}{9E9E9E}         % gris — commentaires
\definecolor{number}{HTML}{E65100}          % orange — nombres
\definecolor{preprocessor}{HTML}{8E24AA}    % violet — preprocesseur
\definecolor{type}{HTML}{00838F}            % teal — types

% ---------- Titres colorés ----------
\usepackage{titlesec}
\titleformat{\section}
  {\normalfont\Large\bfseries\color{sectioncolor}}{\thesection}{1em}{}
\titleformat{\subsection}
  {\normalfont\large\bfseries\color{subsectioncolor}}{\thesubsection}{1em}{}

% Commande pour les questions (violet, en gras, sans numéro)
\newcommand{\question}[1]{%
  \bigskip\noindent{\large\bfseries\color{questioncolor}#1}\par\medskip}

% ---------- Liens / images / tables ----------
\usepackage{hyperref}
\hypersetup{
  colorlinks=true,
  linkcolor=sectioncolor,
  urlcolor=keyword,
  citecolor=sectioncolor
}
\usepackage{graphicx}
\usepackage{array}
\usepackage{longtable}
\usepackage{booktabs}

% ---------- Code avec coloration syntaxique ----------
\usepackage{listings}
\lstdefinestyle{colored}{
  basicstyle=\ttfamily\small,
  columns=fullflexible,
  breaklines=true,
  frame=single,
  framerule=0.5pt,
  rulecolor=\color{codeframe},
  backgroundcolor=\color{codebg},
  showstringspaces=false,
  tabsize=2,
  % Coloration syntaxique
  keywordstyle=\color{keyword}\bfseries,
  stringstyle=\color{string},
  commentstyle=\color{comment}\itshape,
  numberstyle=\tiny\color{comment},
  directivestyle=\color{preprocessor},
  % Numéros de ligne
  numbers=left,
  numbersep=8pt,
  stepnumber=1,
  % Mots-clés supplémentaires
  morekeywords={nullptr,bool,atomic_bool,atomic_flag,assert,mthread_mutex_t,
    mthread_thread_t,mthread_list_item_t,mthread_mutex_attr_t,mthread_sem_t,
    mthread_cond_t,int,void,return,if,else,for,while,struct,typedef},
  emph={mthread_mutex_lock,mthread_mutex_unlock,mthread_mutex_trylock,
    mthread_mutex_init,mthread_mutex_destroy,mthread_self,mthread_create_thread,
    mthread_join,mthread_yield,mthread_block_self,mthread_wake_thread,
    mthread_internal_spin_lock,mthread_internal_spin_unlock,
    atomic_load,atomic_store,atomic_compare_exchange_strong,
    insert_tail_in_list,remove_head_in_list,swapcontext,
    pth_spawn,pth_join,pth_yield,pth_sleep,pth_init,pth_kill,
    pthread_create,pthread_join,printf,fprintf,malloc,free,sleep,
    mthread_init_scheduler_fifo,mthread_log,not_implemented},
  emphstyle=\color{type},
}
\lstset{style=colored}

% ---------- Style pour le code inline ----------
\let\oldtexttt\texttt
\renewcommand{\texttt}[1]{\oldtexttt{\color{keyword}#1}}

\title{\color{sectioncolor}Rapport TD2 — Programmation à base de threads (MPP)}
\author{}
\date{}

\begin{document}
\maketitle

\noindent\textbf{Sujet :} Bibliothèques de threads de niveau utilisateur (GnuPth) et de niveau noyau

\bigskip
\hrule
\bigskip

\section{Partie I — Fonctionnement d'une bibliothèque de niveau utilisateur : GnuPth}

\question{Q.1 : Caractéristiques d'une bibliothèque de threads de niveau utilisateur}

D'après le schéma (processus multithread avec ordonnanceur utilisateur au-dessus de l'ordonnanceur système), les caractéristiques sont :

\begin{itemize}
  \item \textbf{Performances :} Très bonnes pour la création, destruction et changement de contexte entre threads, car tout se fait en espace utilisateur sans appel système. Le coût d'un \emph{context switch} est beaucoup plus faible qu'avec des threads noyau puisqu'on n'a pas besoin de passer en mode noyau.

  \item \textbf{Flexibilité :} Très élevée. L'ordonnanceur est en espace utilisateur, on peut donc le personnaliser et l'adapter à des besoins spécifiques (politique d'ordonnancement, priorités, etc.) sans modifier le noyau du système d'exploitation.

  \item \textbf{SMP (Symmetric MultiProcessing) :} \textbf{Non exploitable}. Puisque l'OS ne voit qu'un seul processus (un seul thread noyau), tous les threads utilisateur s'exécutent sur un seul CPU. L'ordonnanceur utilisateur ne peut pas répartir les threads sur plusieurs cœurs. Il n'y a donc \textbf{aucun parallélisme réel}, seulement de la concurrence (multiplexage temporel).

  \item \textbf{Appels systèmes bloquants :} C'est le \textbf{problème majeur} de ce type de bibliothèque. Si un thread utilisateur effectue un appel système bloquant (ex : \texttt{read()}, \texttt{accept()}), c'est le processus entier qui est bloqué par l'OS, et donc \textbf{tous les autres threads} sont bloqués aussi. L'ordonnanceur utilisateur ne reprend la main qu'au retour de l'appel système. GnuPth contourne partiellement ce problème en remplaçant les appels bloquants par des versions non-bloquantes + \texttt{select()}/\texttt{poll()} via des wrappers.
\end{itemize}

\begin{center}
\begin{tabular}{|l|l|}
\hline
\textbf{Critère} & \textbf{Threads niveau utilisateur (GnuPth)} \\
\hline
Performances (context switch) & \cmark\ Très rapide (pas de syscall) \\
\hline
Flexibilité de l'ordonnanceur & \cmark\ Totale (en espace utilisateur) \\
\hline
Exploitation SMP / multi-cœur & \xmark\ Impossible (1 seul thread noyau) \\
\hline
Appels systèmes bloquants & \xmark\ Bloquent tout le processus \\
\hline
Portabilité & \cmark\ Pas de dépendance au noyau \\
\hline
\end{tabular}
\end{center}

\section{Partie II — Fonctionnement d'une bibliothèque de niveau système : Pthread}

\question{Q.2 : Caractéristiques d'une bibliothèque de threads de niveau système (noyau)}

D'après le schéma (chaque thread utilisateur est associé directement à un thread noyau, géré par l'ordonnanceur système), les caractéristiques sont :

\begin{itemize}
  \item \textbf{Performances :} \textbf{Moins bonnes} que les threads utilisateur pour la création, destruction et le changement de contexte, car chaque opération nécessite un \textbf{appel système} (passage en mode noyau). Le \emph{context switch} est plus coûteux puisqu'il implique une transition \emph{user-space} → \emph{kernel-space} → \emph{user-space}.

  \item \textbf{Flexibilité :} \textbf{Limitée}. L'ordonnancement est entièrement géré par le noyau du système d'exploitation. Le programmeur ne peut pas personnaliser la politique d'ordonnancement (sauf via quelques paramètres comme la priorité ou la politique \texttt{SCHED\_FIFO}/\texttt{SCHED\_RR}). La bibliothèque est aussi \textbf{dépendante du système} sur lequel elle tourne.

  \item \textbf{SMP (Symmetric MultiProcessing) :} \textbf{Pleinement exploitable}. Puisque chaque thread utilisateur correspond à un thread noyau distinct, l'ordonnanceur système peut les répartir sur \textbf{différents CPUs/cœurs}. On obtient donc un \textbf{vrai parallélisme} : plusieurs threads peuvent s'exécuter simultanément sur plusieurs processeurs.

  \item \textbf{Appels systèmes bloquants :} \textbf{Pas de problème}. Si un thread effectue un appel système bloquant, seul \textbf{ce thread} est bloqué. Les autres threads du processus continuent de s'exécuter normalement sur les autres cœurs, car l'ordonnanceur système gère chaque thread indépendamment.
\end{itemize}

\begin{center}
\begin{tabular}{|l|l|}
\hline
\textbf{Critère} & \textbf{Threads niveau système (Pthread)} \\
\hline
Performances (context switch) & \xmark\ Plus lent (appel système nécessaire) \\
\hline
Flexibilité de l'ordonnanceur & \xmark\ Limitée (géré par le noyau) \\
\hline
Exploitation SMP / multi-cœur & \cmark\ Parallélisme réel sur plusieurs CPUs \\
\hline
Appels systèmes bloquants & \cmark\ Ne bloquent que le thread concerné \\
\hline
Portabilité & \xmark\ Dépendante du système d'exploitation \\
\hline
\end{tabular}
\end{center}

\section{Partie III — Évaluation des bibliothèques Pthread et GnuPth sur SMP (processeur multi-cœur)}

\subsection{1. Évaluation des capacités sur architectures SMP}

\question{Q.3 : Moyen de caractériser si une bibliothèque utilise correctement tous les processeurs/cœurs}

Pour vérifier si une bibliothèque de threads exploite bien tous les cœurs disponibles, on peut écrire un \textbf{programme de test CPU-bound} (calcul intensif) qui lance autant de threads que de cœurs, où chaque thread effectue une boucle de calcul lourde (ex : 500\,000\,000 itérations d'un calcul flottant).

\textbf{Méthode concrète :}
\begin{enumerate}
  \item \textbf{Écrire un programme} qui crée \texttt{N} threads (avec \texttt{N} = nombre de cœurs), chacun exécutant un calcul intensif identique.
  \item \textbf{Observer la charge CPU} pendant l'exécution avec un outil comme \texttt{top}, \texttt{htop} ou \texttt{mpstat} :
  \begin{itemize}
    \item Si \textbf{tous les cœurs sont à \(\sim\)100\%} → la bibliothèque exploite bien le SMP.
    \item Si \textbf{un seul cœur est à 100\%} et les autres au repos → la bibliothèque ne supporte pas le SMP.
  \end{itemize}
  \item \textbf{Mesurer le temps d'exécution} avec la commande \texttt{time} :
  \begin{itemize}
    \item \texttt{real} = temps réel écoulé (\emph{wall clock})
    \item \texttt{user} = temps CPU total consommé par tous les threads
    \item Si \(\texttt{user} \approx N \times \texttt{real}\) → \textbf{parallélisme réel} (N cœurs utilisés)
    \item Si \(\texttt{user} \approx \texttt{real}\) → \textbf{pas de parallélisme} (1 seul cœur)
  \end{itemize}
  \item \textbf{Vérifier le \% CPU} dans la sortie de \texttt{time} :
  \begin{itemize}
    \item \(\sim N \times 100\%\) → SMP exploité
    \item \(\sim 100\%\) → un seul cœur utilisé
  \end{itemize}
\end{enumerate}

C'est exactement ce que font nos programmes \texttt{bu\_noyau.c} (Pthread) et \texttt{bu\_utilisateur.c} (GnuPth).

\question{Q.4 : Évaluation des caractéristiques SMP de la bibliothèque GnuPth}

\textbf{Programme utilisé :} \texttt{bu\_utilisateur.c} — chaque thread GnuPth exécute 500\,000\,000 itérations d'un calcul flottant (\texttt{result += i * 0.000001}), avec des \texttt{pth\_yield()} périodiques (tous les 10M d'itérations) car GnuPth est coopératif.

\begin{quote}
\textbf{Fichier source :} \texttt{SMP\_test/bu\_utilisateur.c}
\end{quote}

\textbf{Code source :}
\begin{lstlisting}[language=C]
/*
 * bu_utilisateur.c — Test SMP avec bibliothèque UTILISATEUR (GnuPth)
 *
 * Chaque thread fait le MÊME calcul CPU-intensif que bu_noyau.c.
 * Avec GnuPth, tous les threads partagent UN SEUL thread noyau
 * → pas de parallélisme réel, un seul cœur utilisé.
 *
 * GnuPth est coopératif : un thread garde le CPU jusqu'à ce qu'il
 * appelle pth_yield(). On insère des yield périodiques pour que
 * tous les threads progressent.
 */

#include <stdio.h>
#include <stdlib.h>
#include <unistd.h>
#include <sys/time.h>
#include "pth.h"

#define ITERATIONS 500000000L
#define YIELD_EVERY 10000000L  /* yield périodique pour la coopération */

typedef struct {
    int id;
    double result;
} thread_arg_t;

void *travail_cpu(void *arg)
{
    thread_arg_t *ta = (thread_arg_t *)arg;
    double result = 0.0;

    for (long i = 0; i < ITERATIONS; i++) {
        result += i * 0.000001;
        /* Yield périodique : nécessaire car GnuPth est coopératif.
           Sans yield, un seul thread monopolise le CPU. */
        if (i % YIELD_EVERY == 0)
            pth_yield(NULL);
    }

    ta->result = result;
    return NULL;
}

double get_time(void)
{
    struct timeval tv;
    gettimeofday(&tv, NULL);
    return tv.tv_sec + tv.tv_usec / 1e6;
}

int main(int argc, char **argv)
{
    int nthreads = (argc > 1) ? atoi(argv[1]) : sysconf(_SC_NPROCESSORS_ONLN);

    if (!pth_init()) {
        fprintf(stderr, "Erreur: pth_init() a échoué\n");
        return 1;
    }

    pth_t *threads = malloc(nthreads * sizeof(pth_t));
    thread_arg_t *args = malloc(nthreads * sizeof(thread_arg_t));

    pth_attr_t attr = pth_attr_new();
    pth_attr_set(attr, PTH_ATTR_JOINABLE, TRUE);

    double start = get_time();

    for (int i = 0; i < nthreads; i++) {
        args[i].id = i;
        threads[i] = pth_spawn(attr, travail_cpu, &args[i]);
    }

    pth_attr_destroy(attr);

    for (int i = 0; i < nthreads; i++)
        pth_join(threads[i], NULL);

    double elapsed = get_time() - start;
    printf("Temps réel : %.2f s\n", elapsed);

    pth_kill();
    free(threads);
    free(args);
    return 0;
}
\end{lstlisting}

\textbf{Compilation :}
\begin{lstlisting}[language=bash]
gcc -O2 -I../pth-2.0.7 -L../pth-2.0.7/.libs bu_utilisateur.c -lpth -o bu_utilisateur
\end{lstlisting}

\textbf{Résultats expérimentaux (machine : 8 cœurs) :}
\begin{center}
\begin{tabular}{|c|c|c|c|c|}
\hline
\textbf{Threads} & \textbf{Temps réel (s)} & \textbf{Temps user (s)} & \textbf{\% CPU} & \textbf{Speedup} \\
\hline
1 & 0.60 & 0.60 & 50\% & $\times 1.0$ \\
\hline
2 & 1.27 & 1.25 & 98\% & $\times 0.47$ \\
\hline
4 & 2.33 & 2.32 & 99\% & $\times 0.26$ \\
\hline
8 & 4.72 & 4.70 & 99\% & $\times 0.13$ \\
\hline
\end{tabular}
\end{center}

\textbf{Observations :}
\begin{itemize}
  \item Le \textbf{temps user \(\approx\) temps réel} dans tous les cas → \(\texttt{user}/\texttt{real} \approx 1\) → \textbf{un seul cœur} est utilisé.
  \item Le \textbf{\% CPU \(\approx\) 99--100\%} (jamais \(>100\%\)) → confirme qu'un seul cœur travaille.
  \item Le temps \textbf{augmente linéairement} avec le nombre de threads : 2 threads \(= 2\times\) le temps de 1 thread, 4 threads \(= 4\times\) etc. Il n'y a \textbf{aucun gain}, c'est même \textbf{plus lent} car on fait \(N\) fois plus de travail séquentiellement.
  \item Le speedup est \textbf{inférieur à 1} : ajouter des threads \textbf{dégrade} les performances.
\end{itemize}

\textbf{Conclusion :} GnuPth \textbf{ne supporte pas} le SMP. C'est une bibliothèque de niveau utilisateur : l'OS ne voit qu'\textbf{un seul processus} avec \textbf{un seul thread noyau}. L'ordonnanceur utilisateur de GnuPth multiplexe les threads utilisateur sur ce unique thread noyau de manière coopérative. Puisqu'il n'y a qu'un thread noyau, l'ordonnanceur système ne peut l'affecter qu'à \textbf{un seul CPU} à la fois. Les threads GnuPth s'exécutent de manière \textbf{séquentielle} (concurrence sans parallélisme).

\question{Q.5 : Évaluation des caractéristiques SMP de la bibliothèque Pthread}

\textbf{Programme utilisé :} \texttt{bu\_noyau.c} — même calcul que \texttt{bu\_utilisateur.c} (500\,000\,000 itérations de \texttt{result += i * 0.000001} par thread), mais avec des threads Pthread (niveau noyau).

\begin{quote}
\textbf{Fichier source :} \texttt{SMP\_test/bu\_noyau.c}
\end{quote}

\textbf{Code source :}
\begin{lstlisting}[language=C]
/*
 * bu_noyau.c — Test SMP avec bibliothèque NOYAU (Pthread)
 *
 * Chaque thread fait un calcul CPU-intensif.
 * Avec Pthread, chaque thread = un thread noyau → vrai parallélisme.
 */

#include <stdio.h>
#include <stdlib.h>
#include <pthread.h>
#include <unistd.h>
#include <sys/time.h>

#define ITERATIONS 500000000L

typedef struct {
    int id;
    double result;
} thread_arg_t;

void *travail_cpu(void *arg)
{
    thread_arg_t *ta = (thread_arg_t *)arg;
    double result = 0.0;

    for (long i = 0; i < ITERATIONS; i++)
        result += i * 0.000001;

    ta->result = result;
    return NULL;
}

double get_time(void)
{
    struct timeval tv;
    gettimeofday(&tv, NULL);
    return tv.tv_sec + tv.tv_usec / 1e6;
}

int main(int argc, char **argv)
{
    int nthreads = (argc > 1) ? atoi(argv[1]) : sysconf(_SC_NPROCESSORS_ONLN);

    pthread_t *threads = malloc(nthreads * sizeof(pthread_t));
    thread_arg_t *args = malloc(nthreads * sizeof(thread_arg_t));

    double start = get_time();

    for (int i = 0; i < nthreads; i++) {
        args[i].id = i;
        pthread_create(&threads[i], NULL, travail_cpu, &args[i]);
    }

    for (int i = 0; i < nthreads; i++)
        pthread_join(threads[i], NULL);

    double elapsed = get_time() - start;
    printf("Temps réel : %.2f s\n", elapsed);

    free(threads);
    free(args);
    return 0;
}
\end{lstlisting}

\textbf{Compilation :}
\begin{lstlisting}[language=bash]
gcc -O2 -pthread bu_noyau.c -o bu_noyau
\end{lstlisting}

\textbf{Résultats expérimentaux (machine : 8 cœurs) :}
\begin{center}
\begin{tabular}{|c|c|c|c|c|}
\hline
\textbf{Threads} & \textbf{Temps réel (s)} & \textbf{Temps user (s)} & \textbf{\% CPU} & \textbf{Speedup} \\
\hline
1 & 0.45 & 0.45 & 39\% & $\times 1.0$ \\
\hline
2 & 0.51 & 0.96 & 187\% & $\times 0.88$ \\
\hline
4 & 0.53 & 1.95 & 372\% & $\times 0.85$ \\
\hline
8 & 1.19 & 4.24 & 348\% & $\times 0.38$ \\
\hline
\end{tabular}
\end{center}

\textbf{Observations :}
\begin{itemize}
  \item Le \textbf{temps user \(\gg\) temps réel} dès 2 threads → \(\texttt{user}/\texttt{real} \approx N\) → \textbf{plusieurs cœurs} travaillent en parallèle.
  \item Le \textbf{\% CPU \(\gg\) 100\%} : 187\% pour 2 threads, 372\% pour 4 threads → les threads s'exécutent sur \textbf{plusieurs cœurs simultanément}.
  \item Le \textbf{temps réel reste quasi constant} entre 1, 2 et 4 threads (\(\sim 0.45\)–0.53s) alors que la charge de travail totale est multipliée par \(N\) → \textbf{parallélisme réel effectif}.
  \item À 8 threads, on observe une légère dégradation due à la contention et au fait que certains cœurs (efficacité/performance sur Apple Silicon) n'ont pas la même vitesse.
\end{itemize}

\textbf{Conclusion :} Pthread \textbf{exploite pleinement} le SMP. Chaque thread utilisateur correspond à un \textbf{thread noyau distinct} (modèle 1:1). L'ordonnanceur du système d'exploitation voit ces threads et les \textbf{répartit sur différents cœurs}. On obtient un \textbf{vrai parallélisme matériel} : plusieurs threads s'exécutent réellement en même temps sur plusieurs processeurs.

{Tableau comparatif récapitulatif}

\begin{center}
\begin{tabular}{|l|c|c|}
\hline
\textbf{Critère SMP} & \textbf{GnuPth (niveau utilisateur)} & \textbf{Pthread (niveau système)} \\
\hline
Cœurs utilisés & 1 seul & Tous disponibles \\
\hline
\% CPU avec 4 threads & $\sim$99\% (1 cœur) & $\sim$372\% (4 cœurs) \\
\hline
Temps réel (4 threads) & 2.33s ($\times 4$ vs 1 thread) & 0.53s ($\approx$ même que 1 thread) \\
\hline
Speedup avec N threads & $< 1$ (dégradation) & $\sim 1$ (parallélisme compense) \\
\hline
\texttt{temps user / temps réel} & $\approx 1$ & $\approx N$ \\
\hline
Parallélisme réel & \xmark\ Aucun & \cmark\ Oui \\
\hline
\end{tabular}
\end{center}

\subsection{2. Évaluation des capacités envers les appels bloquants}

\question{Q.6 : Appel système bloquant pour évaluer la gestion des appels bloquants}

L'appel système choisi est \textbf{\texttt{sleep()}} (ou \texttt{pth\_sleep()} pour GnuPth). C'est un appel bloquant classique qui suspend l'exécution du thread appelant pendant un nombre de secondes donné.

\textbf{Principe du test :}
\begin{itemize}
  \item On crée \textbf{N threads}, chacun effectuant un \texttt{sleep(2)} (2 secondes).
  \item Si la bibliothèque \textbf{gère bien} les appels bloquants : les N threads dorment \textbf{en même temps} → temps total \(\approx\) \textbf{2 s} (quel que soit N).
  \item Si la bibliothèque \textbf{gère mal} les appels bloquants : les sleeps s'exécutent \textbf{séquentiellement} → temps total \(\approx N \times 2\) s.
\end{itemize}

On pourrait aussi utiliser \texttt{read()} sur un pipe/socket ou \texttt{accept()} sur un socket, mais \texttt{sleep()} est le plus simple et le plus déterministe pour ce test.

\question{Q.7 : Évaluation des appels bloquants de la bibliothèque GnuPth}

\textbf{Programme utilisé :} \texttt{bloquant\_utilisateur.c} — chaque thread GnuPth effectue un \texttt{pth\_sleep(2)} (wrapper GnuPth de \texttt{sleep()}).

\begin{quote}
\textbf{Fichier source :} \texttt{SMP\_test/bloquant\_utilisateur.c}
\end{quote}

\textbf{Code source :}
\begin{lstlisting}[language=C]
/*
 * bloquant_utilisateur.c — Test appels bloquants avec GnuPth (niveau utilisateur)
 *
 * On crée N threads. Chaque thread fait un pth_sleep() (appel bloquant).
 * GnuPth intercepte les appels bloquants via des wrappers et les
 * remplace par des versions non-bloquantes + select()/poll().
 */

#include <stdio.h>
#include <stdlib.h>
#include <unistd.h>
#include <sys/time.h>
#include "pth.h"

#define SLEEP_TIME 2  /* secondes de blocage par thread */

typedef struct {
    int id;
} thread_arg_t;

double get_time(void)
{
    struct timeval tv;
    gettimeofday(&tv, NULL);
    return tv.tv_sec + tv.tv_usec / 1e6;
}

void *thread_bloquant(void *arg)
{
    thread_arg_t *ta = (thread_arg_t *)arg;
    double start = get_time();

    printf("  [Thread %d] Début pth_sleep(%d)...\n", ta->id, SLEEP_TIME);
    pth_sleep(SLEEP_TIME);  /* appel bloquant GnuPth (wrapper de sleep) */
    double elapsed = get_time() - start;
    printf("  [Thread %d] Réveil après %.2f s\n", ta->id, elapsed);

    return NULL;
}

int main(int argc, char **argv)
{
    int nthreads = (argc > 1) ? atoi(argv[1]) : 4;

    if (!pth_init()) {
        fprintf(stderr, "Erreur: pth_init() a échoué\n");
        return 1;
    }

    pth_t *threads = malloc(nthreads * sizeof(pth_t));
    thread_arg_t *args = malloc(nthreads * sizeof(thread_arg_t));

    pth_attr_t attr = pth_attr_new();
    pth_attr_set(attr, PTH_ATTR_JOINABLE, TRUE);

    double start = get_time();

    for (int i = 0; i < nthreads; i++) {
        args[i].id = i;
        threads[i] = pth_spawn(attr, thread_bloquant, &args[i]);
    }

    pth_attr_destroy(attr);

    for (int i = 0; i < nthreads; i++)
        pth_join(threads[i], NULL);

    double elapsed = get_time() - start;
    printf("Temps total : %.2f s\n", elapsed);

    pth_kill();
    free(threads);
    free(args);
    return 0;
}
\end{lstlisting}

\textbf{Compilation :}
\begin{lstlisting}[language=bash]
gcc -O2 -I../pth-2.0.7 -L../pth-2.0.7/.libs bloquant_utilisateur.c -lpth -o bloquant_utilisateur
\end{lstlisting}

\textbf{Résultats expérimentaux :}
\begin{center}
\begin{tabular}{|c|c|c|c|c|}
\hline
\textbf{Threads} & \textbf{Temps total (s)} & \textbf{Temps attendu (s)} & \textbf{Pire cas (s)} & \textbf{Résultat} \\
\hline
1 & 2.00 & 2 & 2 & \cmark \\
\hline
4 & 2.01 & 2 & 8 & \cmark \\
\hline
8 & 2.00 & 2 & 16 & \cmark \\
\hline
\end{tabular}
\end{center}

\textbf{Observations :}
\begin{itemize}
  \item Le temps total est \textbf{toujours \(\approx 2\) s}, quel que soit le nombre de threads → les \texttt{pth\_sleep()} s'exécutent \textbf{en parallèle}.
  \item Tous les threads démarrent leur sleep en même temps et se réveillent en même temps.
  \item Avec 8 threads, le pire cas serait 16 s (séquentiel), mais on obtient 2 s → \textbf{facteur \(\times 8\)} de gain.
\end{itemize}

\textbf{Explication :} GnuPth gère \textbf{correctement} les appels bloquants grâce à ses \textbf{wrappers}. L'appel \texttt{pth\_sleep()} ne fait \textbf{pas} un vrai \texttt{sleep()} système (qui bloquerait tout le processus). Au lieu de cela, GnuPth :
\begin{enumerate}
  \item Enregistre un \textbf{événement timer} pour le thread appelant
  \item Met le thread dans l'état \textbf{``en attente''}
  \item \textbf{Passe la main} à l'ordonnanceur utilisateur qui peut exécuter un autre thread
  \item Quand le timer expire, le thread est remis dans la file des threads prêts
\end{enumerate}

C'est le \textbf{point fort} de GnuPth : bien que les threads ne s'exécutent que sur un seul cœur, la bibliothèque intercepte les appels bloquants et les remplace par des mécanismes non-bloquants (\texttt{select()}/\texttt{poll()}).

\bigskip
\hrule
\bigskip

\question{Q.8 : Évaluation des appels bloquants de la bibliothèque Pthread}

\textbf{Programme utilisé :} \texttt{bloquant\_noyau.c} — chaque thread Pthread effectue un \texttt{sleep(2)} (appel système standard).

\begin{quote}
\textbf{Fichier source :} \texttt{SMP\_test/bloquant\_noyau.c}
\end{quote}

\textbf{Code source :}
\begin{lstlisting}[language=C]
/*
 * bloquant_noyau.c — Test appels bloquants avec Pthread (niveau noyau)
 *
 * On crée N threads. Chaque thread fait un sleep() (appel bloquant).
 * Avec Pthread, seul le thread appelant est bloqué → les autres continuent.
 */

#include <stdio.h>
#include <stdlib.h>
#include <pthread.h>
#include <unistd.h>
#include <sys/time.h>

#define SLEEP_TIME 2  /* secondes de blocage par thread */

typedef struct {
    int id;
} thread_arg_t;

double get_time(void)
{
    struct timeval tv;
    gettimeofday(&tv, NULL);
    return tv.tv_sec + tv.tv_usec / 1e6;
}

void *thread_bloquant(void *arg)
{
    thread_arg_t *ta = (thread_arg_t *)arg;
    double start = get_time();

    printf("  [Thread %d] Début sleep(%d)...\n", ta->id, SLEEP_TIME);
    sleep(SLEEP_TIME);  /* appel système BLOQUANT */
    double elapsed = get_time() - start;
    printf("  [Thread %d] Réveil après %.2f s\n", ta->id, elapsed);

    return NULL;
}

int main(int argc, char **argv)
{
    int nthreads = (argc > 1) ? atoi(argv[1]) : 4;

    pthread_t *threads = malloc(nthreads * sizeof(pthread_t));
    thread_arg_t *args = malloc(nthreads * sizeof(thread_arg_t));

    double start = get_time();

    for (int i = 0; i < nthreads; i++) {
        args[i].id = i;
        pthread_create(&threads[i], NULL, thread_bloquant, &args[i]);
    }

    for (int i = 0; i < nthreads; i++)
        pthread_join(threads[i], NULL);

    double elapsed = get_time() - start;
    printf("Temps total : %.2f s\n", elapsed);

    free(threads);
    free(args);
    return 0;
}
\end{lstlisting}

\textbf{Compilation :}
\begin{lstlisting}[language=bash]
gcc -O2 -pthread bloquant_noyau.c -o bloquant_noyau
\end{lstlisting}

\textbf{Résultats expérimentaux :}
\begin{center}
\begin{tabular}{|c|c|c|c|c|}
\hline
\textbf{Threads} & \textbf{Temps total (s)} & \textbf{Temps attendu (s)} & \textbf{Pire cas (s)} & \textbf{Résultat} \\
\hline
1 & 2.01 & 2 & 2 & \cmark \\
\hline
4 & 2.01 & 2 & 8 & \cmark \\
\hline
8 & 2.01 & 2 & 16 & \cmark \\
\hline
\end{tabular}
\end{center}

\textbf{Observations :}
\begin{itemize}
  \item Le temps total est \textbf{toujours \(\approx 2\) s}, quel que soit le nombre de threads → les \texttt{sleep()} s'exécutent \textbf{en parallèle}.
  \item Tous les 8 threads démarrent leur sleep, puis se réveillent tous en même temps (\(\sim 2\) s plus tard).
\end{itemize}

\textbf{Explication :} Pthread gère \textbf{naturellement} les appels bloquants puisque chaque thread est un \textbf{thread noyau indépendant}. Quand un thread fait \texttt{sleep()}, le noyau met \textbf{uniquement ce thread} en état \emph{sleeping} ; les autres threads continuent leur exécution normalement. C'est l'avantage fondamental du modèle 1:1 (un thread utilisateur = un thread noyau).

\bigskip
\hrule
\bigskip

{Tableau comparatif — Appels bloquants}

\begin{center}
\renewcommand{\arraystretch}{1.2}
\setlength{\tabcolsep}{4pt}
\small
\begin{tabularx}{\textwidth}{|>{\raggedright\arraybackslash}p{2.8cm}|>{\raggedright\arraybackslash}X|>{\raggedright\arraybackslash}X|}
\hline
\textbf{Critère} & \textbf{GnuPth (niveau utilisateur)} & \textbf{Pthread (niveau système)} \\
\hline
Mécanisme & Wrappers + \texttt{select()}/\texttt{poll()} & Gestion native par le noyau \\
\hline
\texttt{sleep()} avec 4 threads & \(\sim 2\) s \cmark\ (via \texttt{pth\_sleep}) & \(\sim 2\) s \cmark\ (\texttt{sleep} natif) \\
\hline
\texttt{sleep()} avec 8 threads & \(\sim 2\) s \cmark & \(\sim 2\) s \cmark \\
\hline
Appel bloquant natif (\texttt{sleep()} brut) & \xmark\ Bloquerait tout le processus & \cmark\ Ne bloque que le thread \\
\hline
Nécessite des wrappers & \cmark\ Oui (\texttt{pth\_sleep}, \texttt{pth\_read}\ldots) & \xmark\ Non (appels standard) \\
\hline
Limitation & Ne fonctionne que pour les appels wrappés & Aucune limitation \\
\hline
\end{tabularx}
\end{center}

\textbf{Note importante :} GnuPth gère bien les appels bloquants \textbf{uniquement via ses propres wrappers} (\texttt{pth\_sleep()}, \texttt{pth\_read()}, \texttt{pth\_write()}\ldots). Si on utilisait directement \texttt{sleep()} au lieu de \texttt{pth\_sleep()}, \textbf{tout le processus serait bloqué} car l'appel système passerait par le noyau sans que l'ordonnanceur GnuPth puisse intervenir. Pthread n'a pas cette limitation : les appels système standards fonctionnent correctement.

\section{Partie IV — Fonctionnement d'une bibliothèque mixte : mthread}

Les bibliothèques mixtes ont l'avantage de garder un ordonnanceur en espace utilisateur et donc d'être efficaces et flexibles. Toutefois, la nécessité de faire coopérer deux ordonnanceurs rend l'écriture de telles bibliothèques plus difficile. Le schéma illustre ce type de bibliothèque~: plusieurs \textbf{processeurs virtuels (VP)}, chacun porté par un \textbf{thread noyau} (Pthread), et au-dessus un \textbf{ordonnanceur utilisateur} qui multiplexe les threads utilisateur sur ces VP.

\question{Q.9 : Caractéristiques d'une bibliothèque de threads mixte (mthread)}

D'après le schéma (threads utilisateur multiplexés sur plusieurs VP, chaque VP étant un thread noyau, avec un ordonnanceur utilisateur au-dessus de l'ordonnanceur système), les caractéristiques sont~:

\begin{itemize}
  \item \textbf{Performances~:} \textbf{Très bonnes} \cmark. La création, destruction et le changement de contexte entre threads utilisateur se font \textbf{en espace utilisateur} (via \texttt{swapcontext()}), sans appel système. Le coût est comparable à GnuPth et bien inférieur à Pthread. Seule la création des VP (threads Pthread sous-jacents) nécessite des appels système, mais cela n'arrive qu'une fois à l'initialisation.

  \item \textbf{Flexibilité~:} \textbf{Très élevée} \cmark. L'ordonnanceur est en espace utilisateur, donc \textbf{entièrement personnalisable}. On peut implémenter n'importe quelle politique d'ordonnancement (FIFO, priorités, round-robin\ldots) sans modifier le noyau. Dans mthread, l'ordonnanceur est même interchangeable (structure avec pointeurs de fonctions~: \texttt{find\_next\_thread}, \texttt{reschedule\_current}, \texttt{insert\_new\_thread}, \texttt{loadbalancer}).

  \item \textbf{SMP (Symmetric MultiProcessing)~:} \textbf{Exploitable} \cmark. C'est l'\textbf{avantage majeur} par rapport à GnuPth. Puisque chaque VP est un thread noyau distinct, l'ordonnanceur système répartit les VP sur \textbf{différents cœurs CPU}. L'ordonnanceur utilisateur distribue les threads utilisateur entre les VP (via le \texttt{loadbalancer}). On obtient donc un \textbf{vrai parallélisme}~: des threads utilisateur sur VP différents s'exécutent \textbf{simultanément} sur des cœurs différents. Cependant, le parallélisme est \textbf{limité au nombre de VP} (et non au nombre de threads utilisateur).

  \item \textbf{Appels systèmes bloquants~:} \textbf{Partiellement problématiques} $\triangle$. Si un thread utilisateur effectue un appel système bloquant (ex~: \texttt{sleep()}, \texttt{read()}), c'est le \textbf{VP entier} (thread noyau) qui est bloqué par le noyau. Tous les threads utilisateur affectés à ce VP sont alors bloqués aussi. Cependant, les threads sur les \textbf{autres VP} continuent de fonctionner normalement. C'est mieux que GnuPth (où tout est bloqué) mais moins bien que Pthread (où seul le thread appelant est bloqué).
\end{itemize}

\begin{center}
\begin{tabular}{|l|l|}
\hline
\textbf{Critère} & \textbf{Threads mixte (mthread)} \\
\hline
Performances (context switch) & \cmark\ Très rapide (\texttt{swapcontext} en user-space) \\
\hline
Flexibilité de l'ordonnanceur & \cmark\ Totale (ordonnanceur personnalisable en user-space) \\
\hline
Exploitation SMP / multi-cœur & \cmark\ Oui (via les VP = threads noyau sur plusieurs cœurs) \\
\hline
Appels systèmes bloquants & $\triangle$ Bloquent le VP entier (mais pas les autres VP) \\
\hline
Portabilité & \cmark\ Bonne (utilise Pthread + ucontext) \\
\hline
Complexité d'implémentation & \xmark\ Élevée (coordination de 2 ordonnanceurs) \\
\hline
\end{tabular}
\end{center}

\bigskip

{Tableau comparatif des trois types de bibliothèques}

\begin{center}
\renewcommand{\arraystretch}{1.2}
\setlength{\tabcolsep}{4pt}
\small
\begin{tabularx}{\textwidth}{|>{\raggedright\arraybackslash}p{2.8cm}|>{\raggedright\arraybackslash}X|>{\raggedright\arraybackslash}X|>{\raggedright\arraybackslash}X|}
\hline
\textbf{Critère} & \textbf{GnuPth (utilisateur)} & \textbf{Pthread (noyau)} & \textbf{mthread (mixte)} \\
\hline
Context switch & \cmark\ Très rapide & \xmark\ Lent (syscall) & \cmark\ Très rapide \\
\hline
Flexibilité & \cmark\ Totale & \xmark\ Limitée & \cmark\ Totale \\
\hline
SMP / multi-cœur & \xmark\ Impossible & \cmark\ Plein & \cmark\ Oui (via VP) \\
\hline
Appels bloquants & \xmark\ Bloque tout & \cmark\ Un seul thread & $\triangle$ Bloque le VP \\
\hline
Complexité & Simple & Simple & Complexe \\
\hline
Modèle de mapping & N:1 & 1:1 & M:N \\
\hline
\end{tabularx}
\end{center}

\textbf{Explication du modèle M:N~:} mthread implémente un modèle \textbf{M:N} (M threads utilisateur sur N VP/threads noyau, avec $M \geq N$). Cela combine les avantages des deux approches~:
\begin{itemize}
  \item Les \textbf{performances} du modèle N:1 (context switch rapide en user-space)
  \item Le \textbf{parallélisme} du modèle 1:1 (exécution simultanée sur plusieurs cœurs)
\end{itemize}
Le compromis est la \textbf{complexité}~: il faut gérer la répartition des threads utilisateur entre VP (load balancing), la synchronisation entre VP (locks sur les listes \texttt{incoming\_thread}), et les migrations de threads entre VP.

\bigskip
\hrule
\bigskip

\section{Partie V — Découverte du code de la bibliothèque mthread}

La bibliothèque \textbf{mthread} est une bibliothèque de threads de niveau utilisateur écrite en C++. Elle implémente un ordonnanceur coopératif avec support multi-VP (\emph{Virtual Processor}), où chaque VP est porté par un thread Pthread (noyau).

\subsection{1. Découverte du code}

\question{Q.10 : Localisation de l'ordonnanceur dans le code de mthread}

L'ordonnanceur de mthread est réparti dans plusieurs fichiers :

\begin{center}
\begin{tabular}{|l|l|}
\hline
\textbf{Fichier} & \textbf{Rôle} \\
\hline
\texttt{src/mthread\_scheduler\_fifo.h} & Ordonnanceur FIFO (version struct avec pointeurs de fonctions) \\
\hline
\texttt{src/mthread\_scheduler\_fifo\_class.h} & Ordonnanceur FIFO (version classe C++) \\
\hline
\texttt{src/mthread\_scheduler\_empty.h} & Ordonnanceur vide (squelette à implémenter) \\
\hline
\texttt{src/mthread\_scheduler\_helpers.h} & Fonctions utilitaires de l'ordonnanceur \\
\hline
\texttt{src/mthread\_vp.cpp} & Fonctions principales de yield, blocage, réveil \\
\hline
\end{tabular}
\end{center}

Les \textbf{fonctions clés de l'ordonnanceur} sont définies dans la structure \texttt{fifo\_scheduler\_s} (dans \texttt{src/mthread\_scheduler\_fifo.h}) :

\begin{itemize}
  \item \textbf{\texttt{find\_next\_thread(vp)}} — sélectionne le prochain thread à exécuter
  \item \textbf{\texttt{reschedule\_current(vp)}} — réinsère le thread courant en fin de file (round-robin)
  \item \textbf{\texttt{insert\_new\_thread(vp, thread)}} — insère un nouveau thread dans la file des prêts
  \item \textbf{\texttt{loadbalancer(vp, thread)}} — choisit sur quel VP placer un nouveau thread
\end{itemize}

Le \textbf{cœur de l'ordonnancement} se trouve dans la fonction \texttt{internal\_mthread\_yield()} (dans \texttt{src/mthread\_vp.cpp}, \(\sim\) ligne 97).

\question{Q.11 : Fonctionnement de l'ordonnanceur}

L'ordonnanceur mthread est un \textbf{ordonnanceur FIFO coopératif} (non préemptif). Voici son fonctionnement détaillé :

*{1. Structures de données}

Chaque VP (Virtual Processor) possède dans son ordonnanceur :
\begin{itemize}
  \item \texttt{ready\_thread} : liste chaînée des threads prêts à s'exécuter (file FIFO)
  \item \texttt{incoming\_thread} : liste des threads envoyés par d'autres VP (protégée par un lock)
\end{itemize}

\begin{lstlisting}[language=C++]
// src/mthread_scheduler_fifo.h
std::list<mthread_thread_t *> ready_thread;
std::list<mthread_thread_t *> incoming_thread;
SchedLock lock;
\end{lstlisting}

*{2. Sélection du prochain thread — \texttt{find\_next\_thread()}}

\begin{lstlisting}[language=C++]
// src/mthread_scheduler_fifo.h — find_next_thread
// 1. Récupérer les threads entrants (d'autres VP)
vp->scheduler.lock.lock();
if (vp->scheduler.incoming_thread.size() > 0) {
    vp->scheduler.ready_thread.splice(
        vp->scheduler.ready_thread.begin(),
        vp->scheduler.incoming_thread);
}
vp->scheduler.lock.unlock();

// 2. Prendre le premier thread prêt dans ready_thread
while (vp->next_thread == NULL) {
    if (vp->scheduler.ready_thread.size() > 0) {
        vp->next_thread = vp->scheduler.ready_thread.front();
    }
    // 3. Si le thread n'est pas en état "running", le retirer
    //    (zombie → zombie_joinable, blocked → blocked_ready)
}
\end{lstlisting}

*{3. Changement de contexte — \texttt{internal\_mthread\_yield()}}

\begin{lstlisting}[language=C++]
// src/mthread_vp.cpp — internal_mthread_yield (~ligne 97)
// 1. Appeler find_next_thread() pour trouver le prochain thread
vp->scheduler.find_next_thread(vp);

// 2. Si aucun thread prêt, ordonnancer le thread idle
if (vp->next_thread == nullptr)
    vp->next_thread = vp->idle;

// 3. Si le prochain thread est différent du courant, faire un swapcontext
if (vp->current_thread != vp->next_thread) {
    swapcontext(&(tmp_cur->uc), &(tmp_next->uc));
}

// 4. Réinsérer le thread précédent dans la file
vp->scheduler.reschedule_current(vp);
\end{lstlisting}

*{4. Réordonnancement — \texttt{reschedule\_current()}}

Le thread courant (en tête de \texttt{ready\_thread}) est déplacé en \textbf{queue} de la liste, implémentant ainsi une politique \textbf{round-robin FIFO} :

\begin{lstlisting}[language=C++]
// src/mthread_scheduler_fifo.h — reschedule_current
if (vp->scheduler.ready_thread.size() > 1) {
    mthread_move_head_to_tail(vp->scheduler.ready_thread,
                              vp->scheduler.ready_thread);
}
\end{lstlisting}

\textbf{Résumé du cycle :} \texttt{yield()} → \texttt{find\_next\_thread()} → \texttt{swapcontext()} → \texttt{reschedule\_current()} → le thread précédent va en fin de file.

\question{Q.12 : Localisation des fonctions de manipulation des listes}

Les fonctions de manipulation des listes sont dans \textbf{deux fichiers} :

*{1. Listes chaînées simples (pour mutex/synchronisation) — \texttt{src/mthread\_common\_helpers.h}}

\begin{center}
\begin{tabular}{|l|c|l|}
\hline
\textbf{Fonction} & \textbf{Ligne} & \textbf{Description} \\
\hline
\texttt{insert\_tail\_in\_list(item, list)} & \(\sim\) ligne 60 & Insère un élément en \textbf{queue} de la liste \\
\hline
\texttt{remove\_head\_in\_list(list)} & \(\sim\) ligne 79 & Retire et retourne l'élément en \textbf{tête} \\
\hline
\end{tabular}
\end{center}

Ces fonctions manipulent la structure \texttt{mthread\_list\_t} (définie dans \texttt{include/mthread.h}) :
\begin{lstlisting}[language=C++]
typedef struct mthread_list_s {
    volatile mthread_list_item_t *head = nullptr;
    volatile mthread_list_item_t *tail = nullptr;
} mthread_list_t;
\end{lstlisting}

*{2. Listes C++ STL (pour l'ordonnanceur) — \texttt{src/mthread\_scheduler\_helpers.h}}

\begin{center}
\renewcommand{\arraystretch}{1.2}
\setlength{\tabcolsep}{6pt}
\small
\begin{tabularx}{\textwidth}{|>{\raggedright\arraybackslash}p{6cm}|>{\raggedright\arraybackslash}X|}
\hline
\textbf{Fonction} & \textbf{Description} \\
\hline
\texttt{mthread\_move\_head\_to\_tail(from, to)} &
Déplace la tête d’une liste \texttt{std::list} vers la queue (via \texttt{splice}) \\
\hline
\texttt{mthread\_print\_list(l, idle)} &
Affiche le contenu d’une liste pour le debug \\
\hline
\end{tabularx}
\end{center}

L'ordonnanceur utilise des \texttt{std::list} (listes doublement chaînées C++) pour \texttt{ready\_thread} et \texttt{incoming\_thread}, avec les opérations :
\begin{itemize}
  \item \texttt{push\_front()} — insertion en tête
  \item \texttt{pop\_front()} — retrait en tête
  \item \texttt{splice()} — transfert d'éléments entre listes (O(1))
\end{itemize}

\question{Q.13 : Insertion d'un thread dans la liste des threads prêts}

L'insertion d'un nouveau thread dans la file des prêts se fait via la fonction \textbf{\texttt{insert\_new\_thread()}} dans \texttt{src/mthread\_scheduler\_fifo.h} :

\begin{lstlisting}[language=C++]
// src/mthread_scheduler_fifo.h — insert_new_thread
void (*insert_new_thread)(vp, thread) = [](vp, thread) {
    assert(thread->attr.vp != -1);
    if (thread->attr.vp == vp->id) {
        // Cas 1 : Le thread appartient à CE VP
        // → insertion directe en tête de ready_thread
        vp->scheduler.ready_thread.push_front(thread);
    } else {
        // Cas 2 : Le thread appartient à un AUTRE VP
        // → insertion dans incoming_thread du VP distant (avec lock)
        remote_vp->scheduler.lock.lock();
        remote_vp->scheduler.incoming_thread.push_front(thread);
        remote_vp->scheduler.lock.unlock();
    }
};
\end{lstlisting}

\textbf{Deux cas possibles :}
\begin{enumerate}
  \item \textbf{Thread local} (même VP) : insertion directe en \textbf{tête} de \texttt{ready\_thread} avec \texttt{push\_front()}. Le thread sera le prochain à être considéré par \texttt{find\_next\_thread()}.
  \item \textbf{Thread distant} (autre VP) : insertion en \textbf{tête} de \texttt{incoming\_thread} du VP cible, protégée par un \textbf{lock} (car accès concurrent entre VP). Ces threads seront transférés dans \texttt{ready\_thread} lors du prochain appel à \texttt{find\_next\_thread()} via \texttt{splice()}.
\end{enumerate}

\textbf{Note :} L'insertion se fait en \textbf{tête} (\texttt{push\_front}) car le thread courant est toujours le premier élément de \texttt{ready\_thread}. Insérer en tête place le nouveau thread juste devant le thread idle, après le thread courant en queue de parcours.

\question{Q.14 : Comment bloquer un thread}

Le blocage d'un thread se fait via la fonction \textbf{\texttt{internal\_mthread\_block\_self()}} dans \texttt{src/mthread\_vp.cpp} (\(\sim\) ligne 152) :

\begin{lstlisting}[language=C++]
// src/mthread_vp.cpp — internal_mthread_block_self
template <typename Sched>
static inline void internal_mthread_block_self(
    struct mthread_vp_s<Sched> *vp, mthread_mutex_t *lock) {
    // 1. Mettre le thread courant en état "blocked"
    vp->current_thread->status = mthread_blocked;

    // 2. Libérer un éventuel spinlock enregistré
    if (vp->registered_spinlock != nullptr) {
        mthread_internal_spin_unlock(vp->registered_spinlock);
        vp->registered_spinlock = nullptr;
    }

    // 3. Enregistrer le spinlock du mutex pour libération après yield
    vp->registered_spinlock = &(lock->thread_list_spinlock);

    // 4. Faire un yield → l'ordonnanceur va :
    //    a. Voir que le thread est "blocked"
    //    b. Le retirer de ready_thread (status → blocked_ready)
    //    c. Ordonnancer un autre thread
    internal_mthread_yield(vp);
}
\end{lstlisting}

\textbf{Mécanisme complet du blocage (exemple avec un mutex) :}
\begin{enumerate}
  \item \textbf{Le thread appelle \texttt{mthread\_mutex\_lock()}} (\texttt{src/mthread\_mutex.cpp}) :
  \begin{itemize}
    \item Si le mutex est libre → le thread le prend directement
    \item Si le mutex est pris → le thread s'ajoute à la \textbf{liste d'attente du mutex} (\texttt{insert\_tail\_in\_list}) puis appelle \texttt{mthread\_block\_self()}
  \end{itemize}

  \item \textbf{\texttt{mthread\_block\_self()}} :
  \begin{itemize}
    \item Passe le statut du thread à \texttt{mthread\_blocked}
    \item Appelle \texttt{internal\_mthread\_yield()} pour céder le CPU
  \end{itemize}

  \item \textbf{L'ordonnanceur} (dans \texttt{find\_next\_thread()}) :
  \begin{itemize}
    \item Détecte que le thread est \texttt{mthread\_blocked}
    \item Le retire de \texttt{ready\_thread}
    \item Change son statut en \texttt{mthread\_blocked\_ready}
    \item Continue à chercher un autre thread prêt
  \end{itemize}

  \item \textbf{Le réveil} se fait via \texttt{internal\_mthread\_wake\_thread()} (\texttt{src/mthread\_vp.cpp}, \(\sim\) ligne 180) :
  \begin{itemize}
    \item Attend que le thread soit en état \texttt{mthread\_blocked\_ready}
    \item Remet son statut à \texttt{mthread\_running}
    \item Le réinsère dans la file des prêts via \texttt{insert\_new\_thread()}
  \end{itemize}
\end{enumerate}

\textbf{Diagramme des états :}
\begin{lstlisting}[language=text]
running → blocked → blocked_ready → running
                   (par find_next_thread)  (par wake_thread)
\end{lstlisting}

\subsection{2. Mutex}

\question{Q.15 : Utilité des mutex (avec schéma)}

Un \textbf{mutex} (\emph{MUTual EXclusion}) est un mécanisme de synchronisation qui garantit qu'\textbf{un seul thread} à la fois peut accéder à une \textbf{section critique} (ressource partagée).

\textbf{Schéma du fonctionnement d'un mutex :}
\begin{lstlisting}[language=text]
Thread A                    Mutex (libre)                 Thread B
   │                            │                            │
   ├── mutex_lock() ──────────►│ verrouillé (owner=A)       │
   │                            │                            │
   │  ┌─── Section Critique ──┐ │                            │
   │  │  accès à la ressource │ │◄── mutex_lock() ───────── ┤
   │  │  partagée (ex: var++) │ │    → BLOQUÉ               │
   │  └───────────────────────┘ │    (ajouté à la file      │
   │                            │     d'attente du mutex)    │
   ├── mutex_unlock() ────────►│                            │
   │                            │ verrouillé (owner=B)       │
   │                            │ → RÉVEILLÉ ──────────────►├── reprend
   │                            │                            │  ┌─── Section Critique ──┐
   │                            │                            │  │  accès à la ressource │
   │                            │                            │  └───────────────────────┘
   │                            │◄── mutex_unlock() ────────┤
   │                            │ libre                      │
\end{lstlisting}

\textbf{Sans mutex :} deux threads modifiant une même variable en parallèle produisent des \textbf{data races} (résultats incohérents). Par exemple, si deux threads font \texttt{compteur++} en même temps :
\begin{lstlisting}[language=text]
Thread A: lit compteur (=5)         Thread B: lit compteur (=5)
Thread A: calcule 5+1=6             Thread B: calcule 5+1=6
Thread A: écrit compteur=6          Thread B: écrit compteur=6
→ Résultat : compteur=6 au lieu de 7 !
\end{lstlisting}

\textbf{Avec mutex :} un seul thread accède à la variable à la fois → résultat correct.

\textbf{Structure du mutex dans mthread} (définie dans \texttt{include/mthread.h}) :
\begin{lstlisting}[language=C++]
typedef struct mthread_mutex_s {
    bool is_initialized = false;          // mutex initialisé ?
    atomic_bool is_locked = false;        // mutex verrouillé ? (atomique)
    struct mthread_thread_s *owner = nullptr;  // thread propriétaire
    atomic_flag thread_list_spinlock = ATOMIC_FLAG_INIT;  // spinlock de protection
    mthread_list_t thread_list;           // file d'attente des threads bloqués
} mthread_mutex_t;
\end{lstlisting}

\bigskip
\hrule
\bigskip

\subsection{3. Mise en place des mutex dans mthread}

\question{Q.16 : Fonction \texttt{mthread\_mutex\_init} — Localisation et implémentation}

\begin{quote}
\textbf{Fichier source :} \texttt{mthread/src/mthread\_mutex.cpp} (ligne 108)
\end{quote}

\textbf{Rôle :} Initialise un mutex avant sa première utilisation.

\textbf{Code source :}
\begin{lstlisting}[language=C++]
// src/mthread_mutex.cpp — mthread_mutex_init (ligne 108)
int mthread_mutex_init(mthread_mutex_attr_t *attr, mthread_mutex_t *lock) {
    if (lock == nullptr) {
        return MTHREAD_MUTEX_ERROR_NULL;
    }
    if (attr != nullptr) {
        not_implemented();  // les attributs ne sont pas encore supportés
    }
    if (lock->is_initialized) {
        return MTHREAD_MUTEX_ERROR_INIT;  // déjà initialisé
    }
    { lock->is_initialized = true; }
    return 0;
}
\end{lstlisting}

\textbf{Détail de l'implémentation :}
\begin{enumerate}
  \item \textbf{Vérification du pointeur :} si \texttt{lock == nullptr} → retourne \texttt{MTHREAD\_MUTEX\_ERROR\_NULL}.
  \item \textbf{Attributs :} si des attributs sont fournis (\texttt{attr != nullptr}), la fonction appelle \texttt{not\_implemented()} car les attributs ne sont pas encore supportés.
  \item \textbf{Double initialisation :} si le mutex est déjà initialisé (\texttt{is\_initialized == true}) → retourne \texttt{MTHREAD\_MUTEX\_ERROR\_INIT} (protection contre la réinitialisation).
  \item \textbf{Initialisation :} met simplement \texttt{is\_initialized = true}. Les autres champs (\texttt{is\_locked}, \texttt{owner}, \texttt{thread\_list}) gardent leurs valeurs par défaut (grâce aux initialisateurs C++ dans la structure).
\end{enumerate}

\textbf{Note :} L'initialisation est \textbf{minimaliste} car la structure \texttt{mthread\_mutex\_t} a des valeurs par défaut en C++ (\texttt{is\_locked = false}, \texttt{owner = nullptr}, etc.). La fonction ne fait que marquer le mutex comme valide.

\question{Q.17 : Fonction \texttt{mthread\_mutex\_lock} — Localisation et implémentation}

\begin{quote}
\textbf{Fichier source :} \texttt{mthread/src/mthread\_mutex.cpp} (ligne 7)
\end{quote}

\textbf{Rôle :} Verrouille un mutex. Si le mutex est déjà pris, le thread appelant est \textbf{bloqué} jusqu'à ce que le mutex soit libéré.

\textbf{Code source :}
\begin{lstlisting}[language=C++]
// src/mthread_mutex.cpp — mthread_mutex_lock (ligne 7)
int mthread_mutex_lock(mthread_mutex_t *lock) {
    bool unlocked = false;
    if (lock == nullptr) {
        return MTHREAD_MUTEX_ERROR_NULL;
    } else {
        // Auto-initialisation si nécessaire
        if (!lock->is_initialized) {
            int res = mthread_mutex_init(nullptr, lock);
            if (res != 0) return MTHREAD_MUTEX_ERROR_INIT;
        }

        mthread_thread_t *current_thread = mthread_self();

        // Prendre le spinlock pour protéger l'accès au mutex
        mthread_internal_spin_lock(&(lock->thread_list_spinlock));

        if (!lock->is_locked) {
            // CAS 1 : Mutex libre → le prendre
            lock->is_locked = true;
            lock->owner = current_thread;
            // Le spinlock sera libéré implicitement (retour avant yield)
            return 0;
        } else {
            // CAS 2 : Mutex déjà pris → se bloquer
            mthread_list_item_t item;
            item.thread = current_thread;
            // Ajouter le thread à la file d'attente du mutex
            insert_tail_in_list(&item, &lock->thread_list);
            // Se bloquer (libère le spinlock + yield)
            mthread_block_self(lock);
        }
    }
    return 0;
}
\end{lstlisting}

\textbf{Détail de l'implémentation :}
\begin{enumerate}
  \item \textbf{Vérification :} pointeur null + auto-initialisation si nécessaire.
  \item \textbf{Acquisition du spinlock :} \texttt{mthread\_internal\_spin\_lock()} protège l'accès concurrent à la structure du mutex (busy-waiting atomique avec \texttt{atomic\_flag\_test\_and\_set}).
  \item \textbf{Cas 1 — Mutex libre} (\texttt{!lock->is\_locked}) :
  \begin{itemize}
    \item Met \texttt{is\_locked = true}
    \item Enregistre le thread courant comme \texttt{owner}
    \item Retourne immédiatement
  \end{itemize}
  \item \textbf{Cas 2 — Mutex pris} :
  \begin{itemize}
    \item Crée un \texttt{mthread\_list\_item\_t} contenant le thread courant
    \item L'insère en \textbf{queue} de la \texttt{thread\_list} du mutex (FIFO)
    \item Appelle \texttt{mthread\_block\_self(lock)} (statut \texttt{mthread\_blocked} + yield)
    \item À son réveil, le thread reprend ici et retourne 0
  \end{itemize}
\end{enumerate}

\textbf{Remarque :} dans ce code, il y a un \textbf{bug subtil} : dans le cas \emph{mutex libre}, le spinlock n'est pas libéré avant le \texttt{return 0}. Il devrait y avoir un \texttt{mthread\_internal\_spin\_unlock(\&(lock->thread\_list\_spinlock))} juste avant de retourner.

\question{Q.18 : Fonction \texttt{mthread\_mutex\_unlock} — Localisation et implémentation}

\begin{quote}
\textbf{Fichier source :} \texttt{mthread/src/mthread\_mutex.cpp} (ligne 43)
\end{quote}

\textbf{Rôle :} Déverrouille un mutex. S'il y a des threads en attente, réveille le premier de la file.

\textbf{Code source :}
\begin{lstlisting}[language=C++]
// src/mthread_mutex.cpp — mthread_mutex_unlock (ligne 43)
int mthread_mutex_unlock(mthread_mutex_t *lock) {
    bool unlocked = false;
    if (lock == nullptr) {
        return MTHREAD_MUTEX_ERROR_NULL;
    } else {
        // Auto-initialisation si nécessaire
        if (!lock->is_initialized) {
            int res = mthread_mutex_init(nullptr, lock);
            if (res != 0) return MTHREAD_MUTEX_ERROR_INIT;
        }
        // Vérifier que le mutex est bien verrouillé
        if (!atomic_load(&(lock->is_locked))) {
            return MTHREAD_MUTEX_ERROR_NOT_LOCKED;
        }
        // Vérifier que c'est bien le propriétaire qui unlock
        current_thread = mthread_self();
        if (current_thread != lock->owner) {
            return MTHREAD_MUTEX_ERROR_OWNER;
        }

        if (lock->thread_list.head != nullptr) {
            // CAS 1 : Des threads attendent → réveiller le premier
            mthread_internal_spin_lock(&(lock->thread_list_spinlock));
            mthread_list_item_t *item = remove_head_in_list(&(lock->thread_list));
            mthread_internal_spin_unlock(&(lock->thread_list_spinlock));

            // Transférer la propriété au thread réveillé
            lock->owner = item->thread;
            mthread_wake_thread(item->thread);
        } else {
            // CAS 2 : Personne n'attend → libérer le mutex
            atomic_store(&(lock->is_locked), unlocked);  // is_locked = false
        }
    }
    return 0;
}
\end{lstlisting}

\textbf{Détail de l'implémentation :}
\begin{enumerate}
  \item \textbf{Vérifications :}
  \begin{itemize}
    \item Pointeur null → \texttt{MTHREAD\_MUTEX\_ERROR\_NULL}
    \item Mutex non verrouillé → \texttt{MTHREAD\_MUTEX\_ERROR\_NOT\_LOCKED}
    \item Thread appelant \(\neq\) propriétaire → \texttt{MTHREAD\_MUTEX\_ERROR\_OWNER}
  \end{itemize}

  \item \textbf{Cas 1 — Des threads en attente} (\texttt{thread\_list.head != nullptr}) :
  \begin{itemize}
    \item Prend le spinlock pour protéger la file
    \item Retire le \textbf{premier thread} (\texttt{remove\_head\_in\_list} → FIFO)
    \item Libère le spinlock
    \item \textbf{Transfère la propriété} du mutex au thread réveillé (\texttt{lock->owner = item->thread})
    \item \textbf{Réveille} le thread via \texttt{mthread\_wake\_thread()}
    \item Le mutex reste \textbf{verrouillé} : il change de propriétaire
  \end{itemize}

  \item \textbf{Cas 2 — Personne n'attend :}
  \begin{itemize}
    \item Met \texttt{is\_locked = false} (via \texttt{atomic\_store})
    \item Le prochain thread pourra l'acquérir directement
  \end{itemize}
\end{enumerate}

\textbf{Schéma du transfert de propriété lors du unlock :}
\begin{lstlisting}[language=text]
File d'attente du mutex : [Thread B] → [Thread C] → [Thread D]

Thread A appelle unlock() :
  1. Retire Thread B de la file
  2. owner = Thread B (transfert)
  3. wake_thread(B) → B reprend dans lock() et retourne 0

File d'attente restante : [Thread C] → [Thread D]
Mutex toujours verrouillé, owner = Thread B
\end{lstlisting}

\question{Q.19 : Fonction \texttt{mthread\_mutex\_destroy} — Localisation et implémentation}

La fonction \texttt{mthread\_mutex\_destroy} n'existait pas dans le code fourni. Elle n'était ni déclarée dans \texttt{include/mthread.h}, ni implémentée dans \texttt{src/mthread\_mutex.cpp}. Il a fallu la créer.

\begin{quote}
\textbf{Fichier source :} \texttt{mthread/src/mthread\_mutex.cpp} (ajoutée en fin de fichier)\\
\textbf{Déclaration :} \texttt{mthread/include/mthread.h} (ajoutée après \texttt{mthread\_mutex\_init})
\end{quote}

\textbf{Rôle :} Détruire un mutex — opération inverse de \texttt{mthread\_mutex\_init()}. Elle remet la structure dans un état non initialisé, en libérant logiquement la ressource.

\textbf{Code source :}
\begin{lstlisting}[language=C++]
// src/mthread_mutex.cpp — mthread_mutex_destroy
int mthread_mutex_destroy(mthread_mutex_t *lock) {
  if (lock == nullptr) {
    return MTHREAD_MUTEX_ERROR_NULL;
  }
  if (!lock->is_initialized) {
    return MTHREAD_MUTEX_ERROR_INIT;
  }
  if (atomic_load(&(lock->is_locked))) {
    return MTHREAD_MUTEX_ERROR_ALREADY_LOCKED;
  }
  if (lock->thread_list.head != nullptr) {
    return MTHREAD_MUTEX_ERROR_ALREADY_LOCKED;
  }
  lock->is_initialized = false;
  lock->is_locked = false;
  lock->owner = nullptr;
  lock->thread_list.head = nullptr;
  lock->thread_list.tail = nullptr;
  return 0;
}
\end{lstlisting}

\textbf{Détail de l'implémentation :}
\begin{enumerate}
  \item \textbf{Vérification du pointeur NULL :} si \texttt{lock == nullptr} → retourne \texttt{MTHREAD\_MUTEX\_ERROR\_NULL}.
  \item \textbf{Vérification de l'initialisation :} si le mutex n'a jamais été initialisé (\texttt{!is\_initialized}) → retourne \texttt{MTHREAD\_MUTEX\_ERROR\_INIT}.
  \item \textbf{Vérification du verrouillage :} test atomique \texttt{atomic\_load(\&(lock->is\_locked))}. Détruire un mutex verrouillé est un comportement indéfini en POSIX ; ici on retourne une erreur.
  \item \textbf{Vérification de la file d'attente :} si \texttt{thread\_list.head != nullptr}, des threads attendent → on refuse la destruction.
  \item \textbf{Réinitialisation :} remise des champs à leurs valeurs par défaut (\texttt{is\_initialized=false}, \texttt{owner=nullptr}, file vidée).
\end{enumerate}

\textbf{Comparaison avec \texttt{pthread\_mutex\_destroy} (POSIX) :}
\begin{itemize}
  \item POSIX : détruire un mutex verrouillé ou avec des threads en attente → \textbf{comportement indéfini}.
  \item Ici : implémentation \textbf{plus sûre} (codes d'erreur explicites).
\end{itemize}

\question{Q.20 : Fonction \texttt{mthread\_mutex\_trylock} — Implémentation et explication}

\begin{quote}
\textbf{Fichier source :} \texttt{mthread/src/mthread\_mutex.cpp} (ligne 84)
\end{quote}

\textbf{Rôle :} Tentative \textbf{non bloquante} de verrouillage d'un mutex. Si le mutex est libre, il est acquis. S'il est déjà pris, la fonction retourne \textbf{immédiatement} avec un code d'erreur, sans bloquer le thread appelant.

\textbf{Code source :}
\begin{lstlisting}[language=C++]
// src/mthread_mutex.cpp — mthread_mutex_trylock (ligne 84)
int mthread_mutex_trylock(mthread_mutex_t *lock) {
  bool unlocked = false;
  if (lock == nullptr) {
    return MTHREAD_MUTEX_ERROR_NULL;
  } else {
    if (!lock->is_initialized) {
      int res;
      res = mthread_mutex_init(nullptr, lock);
      if (res != 0) {
        return MTHREAD_MUTEX_ERROR_INIT;
      }
    }
    assert(lock != nullptr);
    assert(lock->is_initialized == true);
    if (atomic_compare_exchange_strong(&(lock->is_locked), &unlocked, true)) {
      lock->owner = mthread_self();
      mthread_log("MUTEX", "TryLock %p hold %p\n", lock->owner, lock);
      return 0;
    } else {
      return MTHREAD_MUTEX_ERROR_ALREADY_LOCKED;
    }
  }
  return 0;
}
\end{lstlisting}

\textbf{Détail de l'implémentation :}
\begin{enumerate}
  \item \textbf{Vérifications standard} (comme \texttt{lock}) :
  \begin{itemize}
    \item Pointeur null → \texttt{MTHREAD\_MUTEX\_ERROR\_NULL}
    \item Auto-initialisation si \texttt{!is\_initialized}
  \end{itemize}

  \item \textbf{Tentative atomique avec \texttt{atomic\_compare\_exchange\_strong}} :
  \begin{itemize}
    \item On prépare \texttt{unlocked = false} (valeur attendue dans \texttt{is\_locked})
    \item L'opération CAS fait atomiquement :
    \begin{itemize}
      \item Compare \texttt{lock->is\_locked} avec \texttt{unlocked} (= \texttt{false})
      \item Si \texttt{is\_locked == false} : met \texttt{is\_locked = true} → succès
      \item Si \texttt{is\_locked == true} : ne modifie rien → échec
    \end{itemize}
  \end{itemize}

  \item \textbf{Cas succès} (CAS réussi) :
  \begin{itemize}
    \item Enregistre le thread courant comme propriétaire : \texttt{lock->owner = mthread\_self()}
    \item Retourne \texttt{0}
  \end{itemize}

  \item \textbf{Cas échec} (mutex déjà verrouillé) :
  \begin{itemize}
    \item Retourne \texttt{MTHREAD\_MUTEX\_ERROR\_ALREADY\_LOCKED}
    \item \textbf{Aucun blocage}, aucune insertion dans une file d'attente
  \end{itemize}
\end{enumerate}

\textbf{Différence fondamentale avec \texttt{mthread\_mutex\_lock} :}
\begin{center}
\renewcommand{\arraystretch}{1.2}
\setlength{\tabcolsep}{6pt}
\small
\begin{tabularx}{\textwidth}{|>{\raggedright\arraybackslash}p{3.2cm}|>{\raggedright\arraybackslash}X|>{\raggedright\arraybackslash}X|}
\hline
\textbf{Aspect} & \textbf{lock} & \textbf{trylock} \\
\hline
Mutex libre & Acquis \cmark & Acquis \cmark \\
\hline
Mutex pris & Thread bloqué (yield) & Retourne immédiatement avec erreur \\
\hline
Utilise spinlock & Oui (\texttt{mthread\_internal\_spin\_lock}) & Non (CAS atomique) \\
\hline
Utilise file d’attente & Oui (\texttt{insert\_tail\_in\_list}) & Non \\
\hline
Appelle \texttt{mthread\_block\_self} & Oui & Non \\
\hline
\end{tabularx}
\end{center}

\textbf{Pourquoi \texttt{trylock} n'utilise pas de spinlock ?}  
Parce que \texttt{trylock} ne modifie jamais la \texttt{thread\_list}. Le CAS atomique sur \texttt{is\_locked} suffit pour garantir qu'un seul thread acquiert le mutex à la fois.

\textbf{Cas d'usage typique :}
\begin{lstlisting}[language=C++]
// Polling : essayer d'acquérir le mutex sans bloquer
if (mthread_mutex_trylock(&m) == 0) {
    // Section critique — le mutex est acquis
    shared_resource++;
    mthread_mutex_unlock(&m);
} else {
    // Le mutex est pris → faire autre chose en attendant
    do_other_work();
}
\end{lstlisting}

\bigskip
\hrule
\bigskip

\question{Q.21 : Macro \texttt{MTHREAD\_MUTEX\_INITIALIZER} — Implémentation et explication}

\begin{quote}
\textbf{Fichier :} \texttt{mthread/include/mthread.h} (ajoutée après la définition de \texttt{mthread\_mutex\_t})
\end{quote}

\textbf{Rôle :} Permettre l'initialisation \textbf{statique} d'un mutex, sans appeler \texttt{mthread\_mutex\_init()}.

\textbf{Implémentation :}
\begin{lstlisting}[language=C]
// include/mthread.h — après la définition de mthread_mutex_t
#define MTHREAD_MUTEX_INITIALIZER {       \
    .is_initialized = true,               \
    .is_locked = false,                   \
    .owner = nullptr,                     \
    .thread_list_spinlock = ATOMIC_FLAG_INIT, \
    .thread_list = {nullptr, nullptr}     \
}
\end{lstlisting}

\textbf{Explication détaillée :} c'est un \textbf{designated initializer} C/C++ qui initialise chaque champ de \texttt{mthread\_mutex\_s} :

\begin{center}
\begin{tabular}{|l|l|l|}
\hline
\textbf{Champ} & \textbf{Valeur} & \textbf{Signification} \\
\hline
\texttt{is\_initialized} & \texttt{true} & Mutex prêt à l'emploi \\
\hline
\texttt{is\_locked} & \texttt{false} & Mutex libre \\
\hline
\texttt{owner} & \texttt{nullptr} & Aucun propriétaire \\
\hline
\texttt{thread\_list\_spinlock} & \texttt{ATOMIC\_FLAG\_INIT} & Spinlock initialisé (déverrouillé) \\
\hline
\texttt{thread\_list} & \texttt{\{nullptr, nullptr\}} & File d'attente vide \\
\hline
\end{tabular}
\end{center}

\textbf{Utilisation :}
\begin{lstlisting}[language=C++]
// AVANT (initialisation dynamique) :
mthread_mutex_t m;
mthread_mutex_init(nullptr, &m);

// APRÈS (initialisation statique) :
mthread_mutex_t m = MTHREAD_MUTEX_INITIALIZER;
\end{lstlisting}

\textbf{Avantages de la macro :}
\begin{enumerate}
  \item \textbf{Global/statiques :} permet l'init au niveau fichier (avant \texttt{main}) :
\begin{lstlisting}[language=C++]
// Au niveau fichier (global)
static mthread_mutex_t global_mutex = MTHREAD_MUTEX_INITIALIZER;
\end{lstlisting}
  \item \textbf{Pas de risque d'oubli :} mutex prêt dès la déclaration.
  \item \textbf{Compatibilité POSIX :} équivalent de \texttt{PTHREAD\_MUTEX\_INITIALIZER}.
\end{enumerate}

\textbf{Note :} Avec cette macro, la branche d'auto-initialisation dans \texttt{lock}/\texttt{unlock}/\texttt{trylock} n'est normalement jamais prise car \texttt{is\_initialized} vaut déjà \texttt{true}.

\bigskip
\hrule
\bigskip

\question{Q.22 : Programmes de test pour \texttt{mutex\_destroy}, \texttt{trylock} et \texttt{MTHREAD\_MUTEX\_INITIALIZER}}

\begin{quote}
\textbf{Fichier source :} \texttt{mthread/tests/test\_scheduler\_unit\_mutex\_destroy.cpp}
\end{quote}

\textbf{Programme de test complet :}
\begin{lstlisting}[language=C++]
#include <mthread.h>
#include <cassert>
#include <cstdio>

static int shared_counter = 0;
static mthread_mutex_t mutex_lock;

// ═══════════ Mutex initialisé statiquement (Q.21) ═══════════
static mthread_mutex_t static_mutex = MTHREAD_MUTEX_INITIALIZER;

// ═══════════ TEST Q.19 : mthread_mutex_destroy ═══════════
void test_mutex_destroy() {
    printf("=== Test mthread_mutex_destroy ===\n");

    mthread_mutex_t m;
    int res;

    // Test 1 : Initialiser puis détruire un mutex libre
    res = mthread_mutex_init(nullptr, &m);
    assert(res == 0);
    res = mthread_mutex_destroy(&m);
    assert(res == 0);
    printf("  [OK] destroy sur mutex non verrouillé\n");

    // Test 2 : destroy(NULL) → erreur
    res = mthread_mutex_destroy(nullptr);
    assert(res == MTHREAD_MUTEX_ERROR_NULL);
    printf("  [OK] destroy(NULL) → MTHREAD_MUTEX_ERROR_NULL\n");

    // Test 3 : destroy sur mutex non initialisé → erreur
    mthread_mutex_t m2;
    m2.is_initialized = false;
    res = mthread_mutex_destroy(&m2);
    assert(res == MTHREAD_MUTEX_ERROR_INIT);
    printf("  [OK] destroy sur mutex non initialisé → MTHREAD_MUTEX_ERROR_INIT\n");

    printf("=== Test mthread_mutex_destroy : PASS ===\n\n");
}

// ═══════════ TEST Q.20 : mthread_mutex_trylock ═══════════
static void *trylock_thread_func(void *arg) {
    // Le mutex est déjà verrouillé par le thread principal
    int res = mthread_mutex_trylock(&mutex_lock);
    if (res == 0) {
        shared_counter++;
        printf("  Thread secondaire : trylock réussi, counter = %d\n", shared_counter);
        mthread_mutex_unlock(&mutex_lock);
    } else {
        printf("  Thread secondaire : trylock échoué → %d (ALREADY_LOCKED)\n", res);
        assert(res == MTHREAD_MUTEX_ERROR_ALREADY_LOCKED);
    }
    return nullptr;
}

void test_mutex_trylock() {
    printf("=== Test mthread_mutex_trylock ===\n");

    int res;
    res = mthread_mutex_init(nullptr, &mutex_lock);
    assert(res == 0);

    // Test 1 : trylock sur mutex libre → succès
    res = mthread_mutex_trylock(&mutex_lock);
    assert(res == 0);
    printf("  [OK] trylock sur mutex libre → succès\n");

    // Test 2 : un autre thread tente trylock sur mutex verrouillé → échec
    mthread_thread_t *t1;
    mthread_create_thread(&t1, trylock_thread_func, nullptr);
    mthread_mutex_unlock(&mutex_lock);
    mthread_join(t1, nullptr);
    printf("  [OK] trylock depuis un autre thread\n");

    // Test 3 : trylock(NULL) → erreur
    res = mthread_mutex_trylock(nullptr);
    assert(res == MTHREAD_MUTEX_ERROR_NULL);
    printf("  [OK] trylock(NULL) → MTHREAD_MUTEX_ERROR_NULL\n");

    mthread_mutex_destroy(&mutex_lock);
    printf("=== Test mthread_mutex_trylock : PASS ===\n\n");
}

// ═══════════ TEST Q.21 : MTHREAD_MUTEX_INITIALIZER ═══════════
static void *static_mutex_thread_func(void *arg) {
    mthread_mutex_lock(&static_mutex);
    shared_counter++;
    printf("  Thread : counter = %d (via MTHREAD_MUTEX_INITIALIZER)\n",
           shared_counter);
    mthread_mutex_unlock(&static_mutex);
    return nullptr;
}

void test_mutex_initializer() {
    printf("=== Test MTHREAD_MUTEX_INITIALIZER ===\n");
    shared_counter = 0;

    // Test 1 : lock/unlock direct sans appeler init
    mthread_mutex_lock(&static_mutex);
    shared_counter = 42;
    printf("  [OK] lock/unlock sans init, counter = %d\n", shared_counter);
    mthread_mutex_unlock(&static_mutex);

    // Test 2 : 5 threads concurrents sur le mutex statique
    shared_counter = 0;
    const int NB = 5;
    mthread_thread_t *threads[NB];
    for (int i = 0; i < NB; i++)
        mthread_create_thread(&threads[i], static_mutex_thread_func, nullptr);
    for (int i = 0; i < NB; i++)
        mthread_join(threads[i], nullptr);

    assert(shared_counter == NB);
    printf("  [OK] %d threads → counter = %d (exclusion mutuelle OK)\n",
           NB, shared_counter);

    mthread_mutex_destroy(&static_mutex);
    printf("=== Test MTHREAD_MUTEX_INITIALIZER : PASS ===\n\n");
}

// ═══════════ MAIN ═══════════
int main() {
    mthread_init_scheduler_fifo(1);

    test_mutex_destroy();
    test_mutex_trylock();
    test_mutex_initializer();

    printf("========================================\n");
    printf("  TOUS LES TESTS SONT PASSES !\n");
    printf("========================================\n");
    return 0;
}
\end{lstlisting}

\textbf{Compilation et exécution :}
\begin{lstlisting}[language=bash]
cd mthread/build
cmake ..
make test_scheduler_unit_mutex_destroy
./tests/test_scheduler_unit_mutex_destroy
\end{lstlisting}

\textbf{Sortie attendue :}
\begin{lstlisting}[language=text]
=== Test mthread_mutex_destroy ===
  [OK] destroy sur mutex non verrouillé
  [OK] destroy(NULL) → MTHREAD_MUTEX_ERROR_NULL
  [OK] destroy sur mutex non initialisé → MTHREAD_MUTEX_ERROR_INIT
=== Test mthread_mutex_destroy : PASS ===

=== Test mthread_mutex_trylock ===
  [OK] trylock sur mutex libre → succès
  Thread secondaire : trylock échoué → 5 (ALREADY_LOCKED)
  [OK] trylock depuis un autre thread
  [OK] trylock(NULL) → MTHREAD_MUTEX_ERROR_NULL
=== Test mthread_mutex_trylock : PASS ===

=== Test MTHREAD_MUTEX_INITIALIZER ===
  [OK] lock/unlock sans init, counter = 42
  Thread : counter = 1 (via MTHREAD_MUTEX_INITIALIZER)
  Thread : counter = 2 (via MTHREAD_MUTEX_INITIALIZER)
  Thread : counter = 3 (via MTHREAD_MUTEX_INITIALIZER)
  Thread : counter = 4 (via MTHREAD_MUTEX_INITIALIZER)
  Thread : counter = 5 (via MTHREAD_MUTEX_INITIALIZER)
  [OK] 5 threads → counter = 5 (exclusion mutuelle OK)
=== Test MTHREAD_MUTEX_INITIALIZER : PASS ===

========================================
  TOUS LES TESTS SONT PASSES !
========================================
\end{lstlisting}

\textbf{Résumé des tests :}
\begin{center}
\begin{tabular}{|c|l|l|c|}
\hline
\textbf{Question} & \textbf{Fonction/Macro testée} & \textbf{Tests effectués} & \textbf{Résultat} \\
\hline
Q.19 & \texttt{mthread\_mutex\_destroy} & destroy libre \cmark, NULL \cmark, non-init \cmark & PASS \\
\hline
Q.20 & \texttt{mthread\_mutex\_trylock} & mutex libre \cmark, mutex pris \cmark, NULL \cmark & PASS \\
\hline
Q.21 & \texttt{MTHREAD\_MUTEX\_INITIALIZER} & lock/unlock sans init \cmark, 5 threads concurrents \cmark & PASS \\
\hline
\end{tabular}
\end{center}
\end{document}